% 本文件由账本已入列事件生成；锚点首次呈现仅连接可核的年序与材料限度。
\section{金朝转换中的北方年序}
以下补入按同年并列的政权与地区组织，而不是以一个朝廷的年号吞没其对手。每个可点击事件名保留账本 ID；叙述只将材料明确记录的动作置入年序，并把实施、规模、损失与动机留在可检验的限度内。

\subsection{1119年：宋与金}
《宋史》、《金史》为这一年留下了可回查的年序入口。

\noindent\textbf{宋与金。}\eventanchor{SYD-1119-JIN-001}{丁巳，金人使李善慶來，遣趙有開報聘，至登州而還}。

\subsection{1120年：宋与金}
《宋史》、《金史》为这一年留下了可回查的年序入口。

\noindent\textbf{宋与金。}\eventanchor{SYD-1120-JIN-002}{二月乙亥，遣趙良嗣使金國}；\eventanchor{SYD-1120-JIN-103}{九月壬寅，金人遣勃堇等來}。

\subsection{1121年：宋与金}
《宋史》、《金史》为这一年留下了可回查的年序入口。

\noindent\textbf{宋与金。}\eventanchor{SYD-1121-JIN-003}{丙午，金人再遣曷魯等來}。

\subsection{1122年：宋与金}
《宋史》、《金史》为这一年留下了可回查的年序入口。

\noindent\textbf{宋与金。}\eventanchor{SYD-1122-JIN-004}{己未，金人遣徒孤且烏歇等來議師期}；\eventanchor{SYD-1122-JIN-104}{甲戌，遣趙良嗣報聘于金國}；\eventanchor{SYD-1122-JIN-158}{金人來約夾攻，命童貫為河北、河東路宣撫使，屯兵于邊以應之，且招諭幽、燕}；\eventanchor{SYD-1122-JIN-195}{戊寅，金人遣李靖等來許山前六州}；\eventanchor{SYD-1122-JIN-228}{戊子，遣趙良嗣報聘于金國}。

\subsection{1123年：宋与金}
《宋史》、《金史》为这一年留下了可回查的年序入口。

\noindent\textbf{宋与金。}\eventanchor{SYD-1123-JIN-005}{庚子，童貫、蔡攸入燕，時燕之職官、富民、金帛、子女先為金人盡掠而去}；\eventanchor{SYD-1123-JIN-105}{金主阿骨打殂，弟吳乞買立}；\eventanchor{SYD-1123-JIN-159}{三月乙卯，金人再遣寧術割等來}；\eventanchor{SYD-1123-JIN-196}{十二月乙巳，金人遣高居慶等來賀正旦}；\eventanchor{SYD-1123-JIN-229}{是月，金人取平州，張覺走燕山，金人索之甚急，命王安中縊殺，函其首送之}。

\subsection{1124年：宋与金}
《宋史》、《金史》为这一年留下了可回查的年序入口。

\noindent\textbf{宋与金。}\eventanchor{SYD-1124-JIN-006}{庚子，金人遣富謨弼等以遺留物來獻}；\eventanchor{SYD-1124-JIN-106}{癸卯，金人遣使來告嗣位}；\eventanchor{SYD-1124-JIN-160}{六年春正月乙卯，為金主輟朝}；\eventanchor{SYD-1124-JIN-197}{秋七月戊子，遣許亢宗賀金國嗣位}；\eventanchor{SYD-1124-JIN-230}{戊寅，遣連南夫弔祭金國}。

\subsection{1125年：宋与金}
《宋史》、《金史》为这一年留下了可回查的年序入口。

\noindent\textbf{宋与金。}\eventanchor{SYD-1125-JIN-007}{己酉，中山奏金人斡離不、粘罕分兩道入攻}。

\subsection{1126年：宋与金}
《宋史》、《金史》为这一年留下了可回查的年序入口。

\noindent\textbf{宋与金。}\eventanchor{SYD-1126-JIN-008}{丙辰，妖人郭京用六甲法，盡令守禦人下城，大啟宣化門出攻金人，兵大敗}；\eventanchor{SYD-1126-JIN-107}{丙辰，詔有告奸人妄言金人復至以恐動居民者，賞之}；\eventanchor{SYD-1126-JIN-161}{丁酉，金人陷真定，都鈐轄劉竧死之}；\eventanchor{SYD-1126-JIN-198}{東兵正將占沆與金人戰於交城縣，死之}；\eventanchor{SYD-1126-JIN-231}{二月丁酉朔，命都統制姚平仲將兵夜襲金人軍，不克而奔}。

\subsection{1127年：南宋与金}
《宋史》、《金史》为这一年留下了可回查的年序入口。

\noindent\textbf{南宋与金。}\eventanchor{SYD-1127-JIN-009}{丁酉，金人立張邦昌為楚帝}；\eventanchor{SYD-1127-JIN-108}{二年春正月辛卯朔，命濟王栩、景王杞出賀金軍，金人亦遣使入賀}；\eventanchor{SYD-1127-JIN-162}{庚子，金人來取宗室，開封尹徐秉哲令民結保，毋藏匿}；\eventanchor{SYD-1127-JIN-199}{金人取至軍中，揆抗論，為所殺}；\eventanchor{SYD-1127-JIN-232}{辛未，金人逼上皇召皇后、皇太子入青城}。

\subsection{1128年：南宋与金}
《宋史》、《金史》为这一年留下了可回查的年序入口。

\noindent\textbf{南宋与金。}\eventanchor{SYD-1128-JIN-010}{丙申，復命宇文虛中為資政殿大學士，充金國祈請使}；\eventanchor{SYD-1128-JIN-109}{丙子，金人陷淮寧府，守臣向子韶死之}；\eventanchor{SYD-1128-JIN-163}{丁丑，遣王貺等充金國軍前通問使}；\eventanchor{SYD-1128-JIN-200}{丁未，東京留守統制官薛廣及金人戰於相州，敗死}；\eventanchor{SYD-1128-JIN-233}{二月丙辰，金人再犯東京，宗澤遣統制閻中立等拒之，中立戰死}。

\subsection{1129年：南宋与金}
《宋史》、《金史》为这一年留下了可回查的年序入口。

\noindent\textbf{南宋与金。}\eventanchor{SYD-1129-JIN-011}{丙辰，遣張邵等充金國軍前通問使}；\eventanchor{SYD-1129-JIN-110}{諜報金人治舟師，將由海道窺江、浙，遣韓世忠控守圌山、福山}；\eventanchor{SYD-1129-JIN-164}{丁卯，遣杜時亮使金軍前}；\eventanchor{SYD-1129-JIN-201}{丁巳，金人犯泰州，守臣曾班以城降}；\eventanchor{SYD-1129-JIN-234}{丁酉，遣崔縱使金軍前}。

\subsection{1130年：金与南宋、南宋与金}
《金史》、《宋史》、《三朝北盟会编》、《建炎以来系年要录》为这一年留下了可回查的年序入口。

\noindent\textbf{金与南宋。}\eventanchor{JIN-1130-EVENT-089}{金扶植刘豫政权经营中原}。

\noindent\textbf{南宋与金。}\eventanchor{SYD-1130-JIN-012}{丙辰，金人犯終南縣，經略使鄭恩戰敗，死之}；\eventanchor{SYD-1130-JIN-111}{丙申，詔劉光世節制諸鎮，守禦通、泰州，伺便襲金人過淮}；\eventanchor{SYD-1130-JIN-165}{丙戌，金人自臨安退兵，命劉光世率兵追之}；\eventanchor{SYD-1130-JIN-202}{丁亥，金人陷汴京，權留守上官悟出奔，為盜所殺}；\eventanchor{SYD-1130-JIN-235}{丁卯，金人立劉豫為帝，國號齊}。

\subsection{1131年：南宋与金}
《宋史》、《金史》为这一年留下了可回查的年序入口。

\noindent\textbf{南宋与金。}\eventanchor{SYD-1131-JIN-013}{金人犯和尚原，吳玠擊敗之}；\eventanchor{SYD-1131-JIN-112}{金人犯秦州，吳玠擊敗之}；\eventanchor{SYD-1131-JIN-166}{金人迫興州，張浚退保閬州，以端明殿學士張深為四川制置使，及參議軍事劉子羽趨益昌}；\eventanchor{SYD-1131-JIN-203}{是月，金人攻張榮縮頭湖水砦，榮擊敗之，來告捷，劉光世以榮知泰州}；\eventanchor{SYD-1131-JIN-236}{趙彬及金人合兵圍慶陽府，守臣楊可升擊敗之}。

\subsection{1132年：南宋与金}
《宋史》、《金史》为这一年留下了可回查的年序入口。

\noindent\textbf{南宋与金。}\eventanchor{SYD-1132-JIN-014}{庚子，金人攻方山原，陝西統制楊政援之，金兵引去}；\eventanchor{SYD-1132-JIN-113}{甲寅，金人復自水洛城來攻，楊政等又敗之}；\eventanchor{SYD-1132-JIN-167}{金人侵熙、秦，關師古擊敗之}；\eventanchor{SYD-1132-JIN-204}{金人陷慶陽府，執楊可升，降之}；\eventanchor{SYD-1132-JIN-237}{遣潘致堯等為金國軍前通問使，附茶藥金幣進兩宮}。

\subsection{1133年：南宋与金}
《宋史》、《金史》为这一年留下了可回查的年序入口。

\noindent\textbf{南宋与金。}\eventanchor{SYD-1133-JIN-015}{丙子，王彥復金州，金兵棄均、房去}；\eventanchor{SYD-1133-JIN-114}{己巳，金人遣兵援劉豫，李橫敗走，潁昌復陷}；\eventanchor{SYD-1133-JIN-168}{己酉，金國元帥府遣李永壽、王翊來見}；\eventanchor{SYD-1133-JIN-205}{時金兵深入至金牛鎮，疑有伏，由褒斜谷引兵還興元，吳玠、劉子羽追擊其後，殺獲甚眾}；\eventanchor{SYD-1133-JIN-238}{是月，金人圍方山原，王似命吳玠發兵救之}。

\subsection{1134年：南宋与金}
《宋史》、《金史》为这一年留下了可回查的年序入口。

\noindent\textbf{南宋与金。}\eventanchor{SYD-1134-JIN-016}{丙戌，吳玠敗金兵，復鳳、秦、隴州}；\eventanchor{SYD-1134-JIN-115}{庚戌，賞承州水砦首領徐康等要擊金兵之功，轉官有差，仍蠲承、楚、泰州水砦民兵賦役十年}；\eventanchor{SYD-1134-JIN-169}{癸亥，劉光世遣統制王德擊金人于滁州之桑根，敗之}；\eventanchor{SYD-1134-JIN-206}{己丑，金人攻承州，韓世忠遣將成閔、解元合兵擊於北門，敗之}；\eventanchor{SYD-1134-JIN-239}{己巳，劉光世遣統制王師晟等率兵夜入南壽春府襲金人，敗之，執偽齊知府王靖}。

\subsection{1135年：南宋与金}
《宋史》、《金史》为这一年留下了可回查的年序入口。

\noindent\textbf{南宋与金。}\eventanchor{SYD-1135-JIN-017}{丁未，戒諸軍戰陳毋殺中原民籍充金兵者}；\eventanchor{SYD-1135-JIN-116}{庚戌，張俊遣統領楊忠閔、王進夾擊金人於淮南岸，敗之，降其將程師回、張延壽}；\eventanchor{SYD-1135-JIN-170}{遣何蘚等奉使金國，通問二帝}；\eventanchor{SYD-1135-JIN-207}{是月，金主晟殂，旻之孫亶立}；\eventanchor{SYD-1135-JIN-240}{辛亥，淮東統制崔德明襲敗金兵於盱眙}。

\subsection{1136年：南宋与金}
《宋史》、《金史》为这一年留下了可回查的年序入口。

\noindent\textbf{南宋与金。}\eventanchor{SYD-1136-JIN-018}{是月，劉豫聞親征，告急于金主亶求援，亶不許，豫自起兵三十萬，命子麟趣合肥，侄猊出渦口，引兵分道入寇}；\eventanchor{SYD-1136-JIN-117}{戊戌，韓世忠攻淮陽軍，及金人戰，敗之}；\eventanchor{SYD-1136-JIN-171}{乙卯，韓世忠引兵攻宿遷縣，統制呼延通與金兵戰，敗之，禽其將孛堇牙合}。

\subsection{1137年：南宋与金}
《宋史》、《金史》为这一年留下了可回查的年序入口。

\noindent\textbf{南宋与金。}\eventanchor{SYD-1137-JIN-019}{庚辰，韓世忠引兵渡淮，逆擊金人于劉冷莊，敗之}；\eventanchor{SYD-1137-JIN-118}{庚子，遣王倫等使金國迎奉梓宮}；\eventanchor{SYD-1137-JIN-172}{癸未，王倫等使還，入見，言金國許還梓宮及皇太后，又許還河南諸州}；\eventanchor{SYD-1137-JIN-208}{壬子，統制呼延通、王權等襲擊金人于淮陽軍，敗之，丁巳，詔六參日，輪行在百官一員轉對}。

\subsection{1138年：南宋与金}
《宋史》、《金史》为这一年留下了可回查的年序入口。

\noindent\textbf{南宋与金。}\eventanchor{SYD-1138-JIN-020}{帝曰：「朕嗣守祖宗基業，豈受金人封冊}；\eventanchor{SYD-1138-JIN-119}{丁-\{丑\}-，金國使張通古、蕭哲與王倫皆來}；\eventanchor{SYD-1138-JIN-173}{丁-\{丑\}-，詔：「金國使來，盡割河南、陝西故地，通好于我，許還梓宮及母兄親族，余無需索}；\eventanchor{SYD-1138-JIN-209}{丁未，金國使烏陵思謀、石慶充與王倫等偕來}；\eventanchor{SYD-1138-JIN-241}{辛-\{丑\}-，詔：「金國遣使入境，欲朕屈己就和，命侍從、台諫詳思條奏}。

\subsection{1139年：南宋与金}
《宋史》、《金史》为这一年留下了可回查的年序入口。

\noindent\textbf{南宋与金。}\eventanchor{SYD-1139-JIN-021}{丙辰，金國以撻懶主和割地，疑其二心，殺之}；\eventanchor{SYD-1139-JIN-120}{丙申，金主詔諭河南諸州以割地歸我之意}；\eventanchor{SYD-1139-JIN-174}{王倫見金主于-\{御\}-林子，被拘于河間，遣其副藍公佐先歸}；\eventanchor{SYD-1139-JIN-210}{乙亥，遣前知宿州趙榮、知壽州王威俱還金國}。

\subsection{1140年：南宋与金}
《宋史》、《金史》为这一年留下了可回查的年序入口。

\noindent\textbf{南宋与金。}\eventanchor{SYD-1140-JIN-022}{-\{岳\}-飛領兵援劉錡，與金人戰于蔡州，敗之，復蔡州}；\eventanchor{SYD-1140-JIN-121}{-\{岳\}-飛遣統制郝晸等與金人戰于鄭州北，復鄭州}；\eventanchor{SYD-1140-JIN-175}{丙辰，-\{岳\}-飛將牛皋及金人戰于京西，敗之}；\eventanchor{SYD-1140-JIN-211}{丙戌，金人陷拱州，守臣王慥死之}；\eventanchor{SYD-1140-JIN-242}{川、陝宣撫副使胡世將屢言金人必渝盟，宜爲備}。

\subsection{1141年：南宋与金}
《宋史》、《金史》为这一年留下了可回查的年序入口。

\noindent\textbf{南宋与金。}\eventanchor{SYD-1141-JIN-023}{丙申，遣劉光遠等充金國通問使}；\eventanchor{SYD-1141-JIN-122}{己卯，統制<u>關師古</u>、<u>李橫</u>擊敗金人于<u>巢縣</u>，復之}；\eventanchor{SYD-1141-JIN-176}{壬午，遣魏良臣、王公亮爲金國稟議使}；\eventanchor{SYD-1141-JIN-212}{壬子，璘率姚仲及金人戰于丁劉圈，敗之}；\eventanchor{SYD-1141-JIN-243}{是月，金人陷濠州，邵隆復陝州}。

\subsection{1142年：南宋与金}
《宋史》、《金史》为这一年留下了可回查的年序入口。

\noindent\textbf{南宋与金。}\eventanchor{SYD-1142-JIN-024}{癸巳，金主許歸梓宮及皇太后，遣何鑄等還}；\eventanchor{SYD-1142-JIN-123}{是月，鄭剛中分畫陝西地界，割商、秦之半畀金國，存上津、豐陽、天水三縣及隴西成紀餘地，棄和尚、方山二原，以大散關為界}；\eventanchor{SYD-1142-JIN-177}{辛亥，遣張中孚、中彥還金國}；\eventanchor{SYD-1142-JIN-213}{乙未，遣沈昭遠等賀金主生辰}。

\subsection{1143年：南宋与金}
《宋史》、《金史》为这一年留下了可回查的年序入口。

\noindent\textbf{南宋与金。}\eventanchor{SYD-1143-JIN-025}{己亥，遣鄭朴等使金賀正旦，王師心等賀金主生辰}；\eventanchor{SYD-1143-JIN-124}{戊戌，洪皓至自金國，入見}。

\subsection{1144年：南宋与金}
《宋史》、《金史》为这一年留下了可回查的年序入口。

\noindent\textbf{南宋与金。}\eventanchor{SYD-1144-JIN-026}{甲午，金人來求淮北人之在南者，詔原者聽還}；\eventanchor{SYD-1144-JIN-125}{秋七月戊午，金人殺王倫於河間府}；\eventanchor{SYD-1144-JIN-178}{十四年春正月丁巳，遣羅汝楫等報謝金國}；\eventanchor{SYD-1144-JIN-214}{乙未，遣林保使金賀正旦，宋之才賀金主生辰}。

\subsection{1145年：南宋与金}
《宋史》、《金史》为这一年留下了可回查的年序入口。

\noindent\textbf{南宋与金。}\eventanchor{SYD-1145-JIN-027}{九月辛酉，遣錢周材使金賀正旦，嚴抑賀金主生辰}；\eventanchor{SYD-1145-JIN-126}{三月甲子，遣敷文閣待制周襟、馬觀國、史願、諸將程師回、馬欽、白常皆還金國}。

\subsection{1146年：南宋与金}
《宋史》、《金史》为这一年留下了可回查的年序入口。

\noindent\textbf{南宋与金。}\eventanchor{SYD-1146-JIN-028}{九月甲戌，命何鑄等為金國祈請使，請國族}；\eventanchor{SYD-1146-JIN-127}{壬子，遣邊知白使金賀正旦，周執羔賀金主生辰}。

\subsection{1147年：南宋与金}
《宋史》、《金史》为这一年留下了可回查的年序入口。

\noindent\textbf{南宋与金。}\eventanchor{SYD-1147-JIN-029}{戊申，遣沈該使金賀正旦，詹大方賀金主生辰}。

\subsection{1148年：南宋与金}
《宋史》、《金史》为这一年留下了可回查的年序入口。

\noindent\textbf{南宋与金。}\eventanchor{SYD-1148-JIN-030}{壬申，遣王墨卿使金賀正旦，陳誠之賀金主生辰}。

\subsection{1149年：南宋与金}
《宋史》、《金史》为这一年留下了可回查的年序入口。

\noindent\textbf{南宋与金。}\eventanchor{SYD-1149-JIN-031}{丙寅，秘閣修撰張邵上秦檜在金國代徽宗與粘罕書稿，詔付史館，以邵為徽猷閣待制}。

\subsection{1150年：南宋与金}
《宋史》、《金史》为这一年留下了可回查的年序入口。

\noindent\textbf{南宋与金。}\eventanchor{SYD-1150-JIN-032}{丙戌，遣堯弼等賀金主即位}；\eventanchor{SYD-1150-JIN-128}{辛酉，遣陳誠之使金賀正旦，王〈日严〉賀金主生辰}。

\subsection{1151年：南宋与金}
《宋史》、《金史》为这一年留下了可回查的年序入口。

\noindent\textbf{南宋与金。}\eventanchor{SYD-1151-JIN-033}{甲申，遣陳夔使金賀正旦，陳相賀金主生辰}；\eventanchor{SYD-1151-JIN-129}{壬戌，遣巫伋等為金國祈請使，請歸淵聖皇帝及皇族、增加帝號等事}。

\subsection{1153年：南宋与金}
《宋史》、《金史》为这一年留下了可回查的年序入口。

\noindent\textbf{南宋与金。}\eventanchor{SYD-1153-JIN-034}{戊午，遣吳㮚使金賀正旦，施鉅賀金主生辰}。

\subsection{1154年：南宋与金}
《宋史》、《金史》为这一年留下了可回查的年序入口。

\noindent\textbf{南宋与金。}\eventanchor{SYD-1154-JIN-035}{戊子，遣沈虛中使金賀正旦，張士襄賀金主生辰}。

\subsection{1155年：南宋与金}
《宋史》、《金史》为这一年留下了可回查的年序入口。

\noindent\textbf{南宋与金。}\eventanchor{SYD-1155-JIN-036}{壬午，遣王岷使金賀正旦，鄭柟賀金主生辰}；\eventanchor{SYD-1155-JIN-130}{三月己酉，右司郎中張士襄自金國使還，坐奉使不肅罷官}。

\subsection{1156年：南宋与金}
《宋史》、《金史》为这一年留下了可回查的年序入口。

\noindent\textbf{南宋与金。}\eventanchor{SYD-1156-JIN-037}{庚寅，遣陳誠之等賀金主尊號禮成}；\eventanchor{SYD-1156-JIN-131}{遣李琳使金賀正旦，葛立方賀金主生辰}。

\subsection{1157年：南宋与金}
《宋史》、《金史》为这一年留下了可回查的年序入口。

\noindent\textbf{南宋与金。}\eventanchor{SYD-1157-JIN-038}{辛巳，遣劉章賀金主生辰，丁亥，湯鵬舉罷}。

\subsection{1158年：南宋与金}
《宋史》、《金史》为这一年留下了可回查的年序入口。

\noindent\textbf{南宋与金。}\eventanchor{SYD-1158-JIN-039}{冬十月丁亥朔，遣沈介使金賀正旦，黃中賀金主生辰}。

\subsection{1159年：南宋与金}
《宋史》、《金史》为这一年留下了可回查的年序入口。

\noindent\textbf{南宋与金。}\eventanchor{SYD-1159-JIN-040}{癸卯，遣周麟之等為金國奉表哀謝使}；\eventanchor{SYD-1159-JIN-132}{六月甲辰朔，遣王綸等為金國奉表稱謝使}；\eventanchor{SYD-1159-JIN-179}{十一月丁亥，遣賀允中等為金國遺留國信使}；\eventanchor{SYD-1159-JIN-215}{是月，金國罷沿邊榷場，惟泗州如舊}；\eventanchor{SYD-1159-JIN-244}{乙酉，王綸使還入見，言金國和好無他}。

\subsection{1160年：南宋与金}
《宋史》、《金史》为这一年留下了可回查的年序入口。

\noindent\textbf{南宋与金。}\eventanchor{SYD-1160-JIN-041}{丁未，遣虞允文使金賀正旦，徐度賀金主生辰}；\eventanchor{SYD-1160-JIN-133}{壬申，淮東總管許世安奏，金主亮至汴京，起重兵五十余萬，屯宿、泗州，謀來攻}；\eventanchor{SYD-1160-JIN-180}{壬子，賀允中使還，言金人必叛盟，宜為之備}；\eventanchor{SYD-1160-JIN-216}{戊午，遣葉義問為金國報謝使}。

\subsection{1161年：南宋与金}
《宋史》、《金史》为这一年留下了可回查的年序入口。

\noindent\textbf{南宋与金。}\eventanchor{SYD-1161-JIN-042}{丙辰，金主亮入廬州，王權自昭關遁，金人追至尉子橋，破敵軍統制姚興戰死，權退保和州}；\eventanchor{SYD-1161-JIN-134}{丙子，虞允文督建康諸軍統制官張振、王琪、時俊、戴皋等以舟師拒金主亮於東採石，戰勝，卻之}；\eventanchor{SYD-1161-JIN-181}{丁丑，虞允文遣水軍統制盛新以舟師擊金人于楊林河口，又敗之}；\eventanchor{SYD-1161-JIN-217}{丁未，命宣撫制置司傳檄契丹、西夏、高麗、渤海諸國及河北、河東、陝西、京東、河南諸路，諭出師共討金人}；\eventanchor{SYD-1161-JIN-245}{鄂州統制楊欽以舟師追敗金人于洪澤鎮}。

\subsection{1162年：南宋与金}
《宋史》、《金史》为这一年留下了可回查的年序入口。

\noindent\textbf{南宋与金。}\eventanchor{SYD-1162-JIN-043}{丁巳，遣洪邁等賀金主即位}；\eventanchor{SYD-1162-JIN-135}{庚戌，金人全師來攻，宣敗績，棄去}；\eventanchor{SYD-1162-JIN-182}{庚寅，王剛破金人於海州}；\eventanchor{SYD-1162-JIN-218}{癸丑，金人圍淮寧府，守臣陳亨祖死之}；\eventanchor{SYD-1162-JIN-246}{癸卯，成閔遣統制杜彥救淮甯，擊敗金人于項城縣}。

\subsection{1163年：南宋与金}
《宋史》、《金史》为这一年留下了可回查的年序入口。

\noindent\textbf{南宋与金。}\eventanchor{SYD-1163-JIN-044}{丙午，復宿州，戮金兵數千人}；\eventanchor{SYD-1163-JIN-136}{庚子，遣王子望等爲金國通問使}；\eventanchor{SYD-1163-JIN-183}{甲辰，顯忠及宏淵敗金人於宿州}；\eventanchor{SYD-1163-JIN-219}{金人攻宿州城，顯忠大敗之}；\eventanchor{SYD-1163-JIN-247}{璘已棄德順，道爲金人所邀，將士死者數萬計}。

\subsection{1164年：南宋与金}
《宋史》、《金史》为这一年留下了可回查的年序入口。

\noindent\textbf{南宋与金。}\eventanchor{SYD-1164-JIN-045}{遣洪適等賀金主生辰}；\eventanchor{SYD-1164-JIN-137}{初，金帥以昉等不許四郡，械系之，昉等不屈，金主命歸之}；\eventanchor{SYD-1164-JIN-184}{甲辰，金人犯六合縣，步軍司統制崔皋擊卻之}；\eventanchor{SYD-1164-JIN-220}{壬午，遣魏杞等爲金國通問使}；\eventanchor{SYD-1164-JIN-248}{十一月乙酉，知楚州魏勝與金人戰，死之，州遂陷，濠州亦陷}。

\subsection{1165年：南宋与金}
《宋史》、《金史》为这一年留下了可回查的年序入口。

\noindent\textbf{南宋与金。}\eventanchor{SYD-1165-JIN-046}{癸未，遣<u>王曮</u>等賀金主生辰}。

\subsection{1166年：南宋与金}
《宋史》、《金史》为这一年留下了可回查的年序入口。

\noindent\textbf{南宋与金。}\eventanchor{SYD-1166-JIN-047}{乙亥，遣梁克家等賀金主生辰}。

\subsection{1167年：南宋与金}
《宋史》、《金史》为这一年留下了可回查的年序入口。

\noindent\textbf{南宋与金。}\eventanchor{SYD-1167-JIN-048}{己亥，遣王瀹賀金主生辰}。

\subsection{1168年：南宋与金}
《宋史》、《金史》为这一年留下了可回查的年序入口。

\noindent\textbf{南宋与金。}\eventanchor{SYD-1168-JIN-049}{甲午，禁歸正人藏匿金人者}；\eventanchor{SYD-1168-JIN-138}{十二月丙申，遣胡元質等賀金主生辰}。

\subsection{1169年：南宋与金}
《宋史》、《金史》为这一年留下了可回查的年序入口。

\noindent\textbf{南宋与金。}\eventanchor{SYD-1169-JIN-050}{十二月己-\{丑\}-，遣司馬伋等賀金主生辰}；\eventanchor{SYD-1169-JIN-139}{先是，海州人時旺聚衆數千來請命，旺尋爲金人所獲，其徒渡淮而南者甚衆，故命滋彈壓之}。

\subsection{1170年：南宋与金}
《宋史》、《金史》为这一年留下了可回查的年序入口。

\noindent\textbf{南宋与金。}\eventanchor{SYD-1170-JIN-051}{是月，遣趙雄等賀金主生辰，別函書請更受書之禮}。

\subsection{1171年：南宋与金}
《宋史》、《金史》为这一年留下了可回查的年序入口。

\noindent\textbf{南宋与金。}\eventanchor{SYD-1171-JIN-052}{壬戌，金遣烏林答天錫等來賀會慶節，天錫要帝降榻問金主起居，虞允文請帝還內，命知閣門事王抃諭天錫以明日見，天錫沮退}；\eventanchor{SYD-1171-JIN-140}{十二月丁未，遣翟紱等賀金主生辰}。

\subsection{1172年：南宋与金}
《宋史》、《金史》为这一年留下了可回查的年序入口。

\noindent\textbf{南宋与金。}\eventanchor{SYD-1172-JIN-053}{丁巳，遣韓元吉等賀金主生辰}；\eventanchor{SYD-1172-JIN-141}{姚憲、曾覿至自金，金人拒其請}。

\subsection{1173年：南宋与金}
《宋史》、《金史》为这一年留下了可回查的年序入口。

\noindent\textbf{南宋与金。}\eventanchor{SYD-1173-JIN-054}{遣韓彥直等賀金主生辰}。

\subsection{1174年：南宋与金}
《宋史》、《金史》为这一年留下了可回查的年序入口。

\noindent\textbf{南宋与金。}\eventanchor{SYD-1174-JIN-055}{壬戌，遣吳琚等賀金主生辰}。

\subsection{1175年：南宋与金}
《宋史》、《金史》为这一年留下了可回查的年序入口。

\noindent\textbf{南宋与金。}\eventanchor{SYD-1175-JIN-056}{遣張宗元等賀金主生辰}。

\subsection{1176年：南宋与金}
《宋史》、《金史》为这一年留下了可回查的年序入口。

\noindent\textbf{南宋与金。}\eventanchor{SYD-1176-JIN-057}{庚午，遣張子正等賀金主生辰}。

\subsection{1177年：南宋与金}
《宋史》、《金史》为这一年留下了可回查的年序入口。

\noindent\textbf{南宋与金。}\eventanchor{SYD-1177-JIN-058}{癸亥，遣趙思等賀金主生辰}。

\subsection{1178年：南宋与金}
《宋史》、《金史》为这一年留下了可回查的年序入口。

\noindent\textbf{南宋与金。}\eventanchor{SYD-1178-JIN-059}{辛卯，遣錢沖之等賀金主生辰}。

\subsection{1179年：南宋与金}
《宋史》、《金史》为这一年留下了可回查的年序入口。

\noindent\textbf{南宋与金。}\eventanchor{SYD-1179-JIN-060}{癸未，遣傅淇等賀金主生辰}。

\subsection{1181年：南宋与金}
《宋史》、《金史》为这一年留下了可回查的年序入口。

\noindent\textbf{南宋与金。}\eventanchor{SYD-1181-JIN-061}{丁酉，遣燕世良賀金主生辰}。

\subsection{1182年：南宋与金}
《宋史》、《金史》为这一年留下了可回查的年序入口。

\noindent\textbf{南宋与金。}\eventanchor{SYD-1182-JIN-062}{丙戌，遣賈選等賀金主生辰}。

\subsection{1183年：南宋与金}
《宋史》、《金史》为这一年留下了可回查的年序入口。

\noindent\textbf{南宋与金。}\eventanchor{SYD-1183-JIN-063}{丁巳，遣陳居仁等賀金主生辰}。

\subsection{1184年：南宋与金}
《宋史》、《金史》为这一年留下了可回查的年序入口。

\noindent\textbf{南宋与金。}\eventanchor{SYD-1184-JIN-064}{癸巳，命利路三都統吳挺、郭鈞、彭杲密陳出師進取利害，以備金人}；\eventanchor{SYD-1184-JIN-142}{盱眙軍言得金人牒，以上京地寒，來歲正旦、生辰人使權止一年}；\eventanchor{SYD-1184-JIN-185}{乙丑，遣王信等賀金主生辰}。

\subsection{1185年：南宋与金}
《宋史》、《金史》为这一年留下了可回查的年序入口。

\noindent\textbf{南宋与金。}\eventanchor{SYD-1185-JIN-065}{壬辰，遣章森等賀金主生辰}。

\subsection{1186年：南宋与金}
《宋史》、《金史》为这一年留下了可回查的年序入口。

\noindent\textbf{南宋与金。}\eventanchor{SYD-1186-JIN-066}{庚申，遣張叔椿等賀金主生辰}。

\subsection{1187年：南宋与金}
《宋史》、《金史》为这一年留下了可回查的年序入口。

\noindent\textbf{南宋与金。}\eventanchor{SYD-1187-JIN-067}{遣胡晉臣等賀金主生辰}；\eventanchor{SYD-1187-JIN-143}{以顏師魯等充金國遣留國信使}。

\subsection{1189年：南宋与金}
《宋史》、《金史》为这一年留下了可回查的年序入口。

\noindent\textbf{南宋与金。}\eventanchor{SYD-1189-JIN-068}{十六年春正月癸巳，金主雍殂，孫璟立}；\eventanchor{SYD-1189-JIN-144}{戊辰，遣謝深甫等賀金主生辰}。

\subsection{1190年：南宋与金}
《宋史》、《金史》为这一年留下了可回查的年序入口。

\noindent\textbf{南宋与金。}\eventanchor{SYD-1190-JIN-069}{六月丁亥，遣丘崈等賀金主生辰}。

\subsection{1191年：南宋与金}
《宋史》、《金史》为这一年留下了可回查的年序入口。

\noindent\textbf{南宋与金。}\eventanchor{SYD-1191-JIN-070}{庚辰，遣趙廱等賀金主生辰}。

\subsection{1192年：南宋与金}
《宋史》、《金史》为这一年留下了可回查的年序入口。

\noindent\textbf{南宋与金。}\eventanchor{SYD-1192-JIN-071}{以權禮部尚書陳騤同知樞密院事，甲辰，遣錢之望等賀金主生辰}。

\subsection{1193年：南宋与金}
《宋史》、《金史》为这一年留下了可回查的年序入口。

\noindent\textbf{南宋与金。}\eventanchor{SYD-1193-JIN-072}{己亥，遣許及之等賀金主生辰}。

\subsection{1194年：南宋与金}
《宋史》、《金史》为这一年留下了可回查的年序入口。

\noindent\textbf{南宋与金。}\eventanchor{SYD-1194-JIN-073}{六月，遣梁總等賀金主生辰}。

\subsection{1195年：南宋与金}
《宋史》、《金史》为这一年留下了可回查的年序入口。

\noindent\textbf{南宋与金。}\eventanchor{SYD-1195-JIN-074}{己未，遣汪義端賀金主生辰}。

\subsection{1196年：南宋与金}
《宋史》、《金史》为这一年留下了可回查的年序入口。

\noindent\textbf{南宋与金。}\eventanchor{SYD-1196-JIN-075}{六月庚戌，遣吳宗旦賀金主生辰，乙丑，命監司、帥守臧否縣令，分三等，丙子，子埈生}。

\subsection{1197年：南宋与金}
《宋史》、《金史》为这一年留下了可回查的年序入口。

\noindent\textbf{南宋与金。}\eventanchor{SYD-1197-JIN-076}{乙亥，遣衛涇賀金主生辰}。

\subsection{1198年：南宋与金}
《宋史》、《金史》为这一年留下了可回查的年序入口。

\noindent\textbf{南宋与金。}\eventanchor{SYD-1198-JIN-077}{六月己巳，遣楊王休賀金主生辰}。

\subsection{1199年：南宋与金}
《宋史》、《金史》为这一年留下了可回查的年序入口。

\noindent\textbf{南宋与金。}\eventanchor{SYD-1199-JIN-078}{六月癸亥，遣李大性賀金主生辰}。

\subsection{1200年：南宋与金}
《宋史》、《金史》为这一年留下了可回查的年序入口。

\noindent\textbf{南宋与金。}\eventanchor{SYD-1200-JIN-079}{丙子，遣丁常任爲金國遺留國信使}；\eventanchor{SYD-1200-JIN-145}{壬辰，遣趙善義賀金主生辰，吳旴使金告哀}。

\subsection{1201年：南宋与金}
《宋史》、《金史》为这一年留下了可回查的年序入口。

\noindent\textbf{南宋与金。}\eventanchor{SYD-1201-JIN-080}{六月辛巳，遣陳宗召賀金主生辰}。

\subsection{1202年：南宋与金}
《宋史》、《金史》为这一年留下了可回查的年序入口。

\noindent\textbf{南宋与金。}\eventanchor{SYD-1202-JIN-081}{六月丙子，遣趙不艱賀金主生辰}。

\subsection{1203年：南宋与金}
《宋史》、《金史》为这一年留下了可回查的年序入口。

\noindent\textbf{南宋与金。}\eventanchor{SYD-1203-JIN-082}{六月壬寅，遣劉甲賀金主生辰}；\eventanchor{SYD-1203-JIN-146}{是冬，金國多難，懼朝廷乘其隙，沿邊聚糧增戍，且禁襄陽榷場}。

\subsection{1204年：南宋与金}
《宋史》、《金史》为这一年留下了可回查的年序入口。

\noindent\textbf{南宋与金。}\eventanchor{SYD-1204-JIN-083}{六月癸巳，遣張嗣古賀金主生辰}。

\subsection{1205年：南宋与金}
《宋史》、《金史》为这一年留下了可回查的年序入口。

\noindent\textbf{南宋与金。}\eventanchor{SYD-1205-JIN-084}{己亥，遣李壁賀金主生辰}；\eventanchor{SYD-1205-JIN-147}{是月，金國以邊民侵掠及增邊戍來責渝盟}。

\subsection{1206年：南宋与金}
《宋史》、《金史》为这一年留下了可回查的年序入口。

\noindent\textbf{南宋与金。}\eventanchor{SYD-1206-JIN-085}{庚戌，金人破成州，守臣辛槱之遁去}；\eventanchor{SYD-1206-JIN-148}{癸卯，郭倬等還至蘄縣，金人追而圍之，倬執馬軍司統制田俊邁以與金人，乃得免}；\eventanchor{SYD-1206-JIN-186}{甲寅，金人攻六合縣，郭倪遣前軍統制郭僎救之，遇於胥浦橋，大敗，倪棄揚州走}；\eventanchor{SYD-1206-JIN-221}{九月壬午，金兵攻奪和尚原}；\eventanchor{SYD-1206-JIN-249}{是月，濠州、安豐軍及邊屯皆爲金人所破}。

\subsection{1207年：南宋与金}
《宋史》、《金史》为这一年留下了可回查的年序入口。

\noindent\textbf{南宋与金。}\eventanchor{SYD-1207-JIN-086}{癸-\{丑\}-，金人復破隨州}；\eventanchor{SYD-1207-JIN-149}{癸亥，以林拱辰爲金國通謝使，遣富琯使金告哀，劉彌正賀金主生辰}；\eventanchor{SYD-1207-JIN-187}{癸酉，金人復破大散關}；\eventanchor{SYD-1207-JIN-222}{己未，奉使金國通謝、國信所參議官方信孺發行在}；\eventanchor{SYD-1207-JIN-250}{四川宣撫副使司參贊軍事楊巨源與金人戰於長橋，敗績}。

\subsection{1208年：南宋与金}
《宋史》、《金史》为这一年留下了可回查的年序入口。

\noindent\textbf{南宋与金。}\eventanchor{SYD-1208-JIN-087}{丙子，遣鄒應龍賀金主生辰}。

\subsection{1209年：南宋与金}
《宋史》、《金史》为这一年留下了可回查的年序入口。

\noindent\textbf{南宋与金。}\eventanchor{SYD-1209-JIN-088}{己巳，遣俞應符賀金主生辰}。

\subsection{1210年：南宋与金}
《宋史》、《金史》为这一年留下了可回查的年序入口。

\noindent\textbf{南宋与金。}\eventanchor{SYD-1210-JIN-089}{癸亥，遣黃中賀金主生辰}。

\subsection{1211年：南宋与金}
《宋史》、《金史》为这一年留下了可回查的年序入口。

\noindent\textbf{南宋与金。}\eventanchor{SYD-1211-JIN-090}{冬十月甲辰，以金國有難，命江淮、京湖、四川制置司謹邊備}；\eventanchor{SYD-1211-JIN-150}{六月丁亥，遣余嶸賀金主生辰，會金國有難，不至而還}。

\subsection{1212年：南宋与金}
《宋史》、《金史》为这一年留下了可回查的年序入口。

\noindent\textbf{南宋与金。}\eventanchor{SYD-1212-JIN-091}{六月癸未，遣傅誠賀金主生辰}。

\subsection{1213年：南宋与金}
《宋史》、《金史》为这一年留下了可回查的年序入口。

\noindent\textbf{南宋与金。}\eventanchor{SYD-1213-JIN-092}{丙戌，以金主新立，命四川謹邊備}；\eventanchor{SYD-1213-JIN-151}{丁-\{丑\}-，遣董居誼賀金主生辰，會金國亂，不至而還}；\eventanchor{SYD-1213-JIN-188}{戊申，遣真德秀賀金主即位，會金國亂，不至而還}。

\subsection{1214年：南宋与金}
《宋史》、《金史》为这一年留下了可回查的年序入口。

\noindent\textbf{南宋与金。}\eventanchor{SYD-1214-JIN-093}{庚辰，金國來督二年歲幣}；\eventanchor{SYD-1214-JIN-152}{庚寅，以起居舍人真德秀奏，罷金國歲幣}；\eventanchor{SYD-1214-JIN-189}{七年春正月丁卯朔，四川制置司遣提舉皂郊博馬務何九齡率諸將及金人戰于秦州城下，敗還}；\eventanchor{SYD-1214-JIN-223}{是月，夏人以書來四川，議夾攻金人，不報}。

\subsection{1215年：南宋与金}
《宋史》、《金史》为这一年留下了可回查的年序入口。

\noindent\textbf{南宋与金。}\eventanchor{SYD-1215-JIN-094}{己卯，遣丁焴賀金主生辰}。

\subsection{1216年：南宋与金}
《宋史》、《金史》为这一年留下了可回查的年序入口。

\noindent\textbf{南宋与金。}\eventanchor{SYD-1216-JIN-095}{乙亥，遣留筠賀金主生辰}。

\subsection{1217年：南宋与金}
《宋史》、《金史》为这一年留下了可回查的年序入口。

\noindent\textbf{南宋与金。}\eventanchor{SYD-1217-JIN-096}{庚子，遣錢撫賀金主生辰}；\eventanchor{SYD-1217-JIN-153}{夏四月丁未朔，金人犯光州中渡鎮，執榷場官盛允升殺之，遂分兵犯樊城}；\eventanchor{SYD-1217-JIN-190}{辛酉，廬州鈐轄王辛敗金人於光山縣之安昌砦，殺其統軍完顏掩}。

\subsection{1218年：南宋与金}
《宋史》、《金史》为这一年留下了可回查的年序入口。

\noindent\textbf{南宋与金。}\eventanchor{SYD-1218-JIN-097}{丙午，金人破皂郊，死者五萬人}；\eventanchor{SYD-1218-JIN-154}{楚州鈐轄梁昭祖焚金人糧舟於大清河，京東忠義副都統沈鐸遣兵助之}；\eventanchor{SYD-1218-JIN-191}{十一月壬申，金人攻安豐軍之黃口灘}；\eventanchor{SYD-1218-JIN-224}{戊午，金人復犯大散關，守將王立遁}。

\subsection{1219年：南宋与金}
《宋史》、《金史》为这一年留下了可回查的年序入口。

\noindent\textbf{南宋与金。}\eventanchor{SYD-1219-JIN-098}{二月戊戌朔，金人破光山縣}；\eventanchor{SYD-1219-JIN-155}{庚寅，金人犯隨州、棗陽軍，又破信陽軍之二砦，京西諸將引兵拒之}；\eventanchor{SYD-1219-JIN-192}{癸卯，金人破武休關，興元都統李貴遁還，利州路提刑、權興元府事趙希昔棄城去}；\eventanchor{SYD-1219-JIN-225}{癸巳，丁焴復以書約夏國攻金人}。

\subsection{1220年：南宋与金}
《宋史》、《金史》为这一年留下了可回查的年序入口。

\noindent\textbf{南宋与金。}\eventanchor{SYD-1220-JIN-099}{庚戌，金人犯皂郊堡，沔州統制董炤等與戰，大敗}；\eventanchor{SYD-1220-JIN-156}{金人追之，遂攻樊城，趙方督諸將拒退之}；\eventanchor{SYD-1220-JIN-193}{壬申，安丙遺夏人書，定議夾攻金人}；\eventanchor{SYD-1220-JIN-226}{壬寅，質俊等自來遠鎮進攻定邊城，金人來救，俊等擊破之}。

\subsection{1221年：南宋与金}
《宋史》、《金史》为这一年留下了可回查的年序入口。

\noindent\textbf{南宋与金。}\eventanchor{SYD-1221-JIN-100}{丁亥，金人破黃州，淮西提刑、知州事何大節棄城遁死}；\eventanchor{SYD-1221-JIN-157}{癸-\{丑\}-，金人退師，扈再興邀擊，敗之於天長鎮，甲寅晦，又敗之}；\eventanchor{SYD-1221-JIN-194}{己亥，金人陷蘄州，知州事李誠之及其家人、官屬皆死之}；\eventanchor{SYD-1221-JIN-227}{壬申，金人治舟於團風，弗克濟，遂圍黃州，分兵破諸縣，又遣別將犯漢陽軍}。

\subsection{1224年：南宋与金}
《宋史》、《金史》为这一年留下了可回查的年序入口。

\noindent\textbf{南宋与金。}\eventanchor{SYD-1224-JIN-101}{癸酉，知西和州尚震午坐金兵至謀遁，奪三官、岳州居住}。

\subsection{1232年：南宋与金}
《宋史》、《金史》为这一年留下了可回查的年序入口。

\noindent\textbf{南宋与金。}\eventanchor{SYD-1232-JIN-102}{時宋與大元兵合圍汴京，金主奔歸德府，尋奔蔡州，大元再遣使議攻金，史嵩之以鄒伸之報謝}。
